% Generated draft: requires mathematical and source review.
\documentclass[11pt]{article}
\usepackage[margin=25mm]{geometry}
\usepackage{fontspec}
\setmainfont{Latin Modern Roman}
\usepackage{amsmath,amssymb,mathtools,unicode-math}
\setmathfont{Latin Modern Math}
\usepackage{enumitem,needspace,xurl,hyperref}
\setlength{\parindent}{0pt}
\setlength{\parskip}{5pt}
\setlength{\emergencystretch}{3em}
\newcommand{\N}{\mathbb{N}}
\newcommand{\Z}{\mathbb{Z}}
\newcommand{\Q}{\mathbb{Q}}
\newcommand{\R}{\mathbb{R}}
\newcommand{\C}{\mathbb{C}}
\newcommand{\F}{\mathbb{F}}
\newcommand{\A}{\mathbb{A}}
\newcommand{\PP}{\mathbb{P}}
\newcommand{\E}{\mathbb{E}}
\newcommand{\eps}{\varepsilon}
\newcommand{\dd}{\,\mathrm{d}}
\begin{document}
\section*{UnsolvedMath ID: 30002376}
\textbf{Source problem code:} catalogue entry
\subsection{Definitions and mathematical statement}
Let \(M\) be a closed, connected, smooth three-dimensional manifold; here closed means compact and without boundary. Its diffeomorphism type is its equivalence class under smooth diffeomorphisms.

A standard interpretation of a contact-type energy-level realization is the following. Let \((W,\omega)\) be a symplectic four-dimensional manifold, let \(H\colon W\to\R\) be smooth, and identify \(M\) with a connected component of a regular level \(H^{-1}(c)\). Suppose there is a vector field \(Y\), defined near \(M\), transverse to \(M\), and satisfying \(\mathcal L_Y\omega=\omega\). For the inclusion \(\iota\colon M\hookrightarrow W\), the one-form
\[
\alpha=\iota^*(\omega(Y,\cdot))
\]
is a contact form: \(\alpha\wedge\dd\alpha\) is nowhere zero. The induced contact structure is the plane field \(\xi=\ker\alpha\subset TM\). The Reeb vector field \(R_\alpha\) is uniquely specified by
\[
\alpha(R_\alpha)=1,
\qquad
\dd\alpha(R_\alpha,\cdot)=0.
\]
Its trajectories follow the characteristic line field \(\ker(\iota^*\omega)\). The Hamiltonian flow restricted to the regular energy level also follows this line field, although its parametrization need not equal the Reeb parametrization. These definitions explain the usual contact-type setting; the supplied record does not specify the ambient realization or identify which parametrized flow is intended.

The supplied problem asks: can the diffeomorphism type of a closed connected three-manifold, regarded as a contact-type energy level, be determined by dynamical properties of the flow on that level together with information about its induced contact structure?

This is a characterization question, not a specified theorem or classification criterion. The record does not prescribe which dynamical properties are available, how contact structures are to be compared, or what constitutes determination of the diffeomorphism type. Consequently, no particular collection of invariants or necessary-and-sufficient criterion can be stated from the supplied evidence alone. Interpreting the input as the entire contact manifold up to contactomorphism would already retain its underlying diffeomorphism type; a substantive interpretation therefore requires a specification of the dynamical and contact information that may be used. The source does not provide that specification.
\subsection{Short English statement}
The question asks whether flow dynamics and induced contact-structure information on a closed connected three-dimensional contact-type energy level can characterize its diffeomorphism type. The admissible data and the intended notion of characterization remain unspecified. The supplied dated curation labels the problem partially solved and reports results for standard tight lens spaces, including the three-sphere; this draft does not independently verify that assessment.
\subsection{Sources}
\begin{itemize}
\item Local UnsolvedMath registry, supplied complete source record: permanent ID \(30002376\), exact source problem code \texttt{catalogue entry}, title ``Recovering Three-Manifold Topology from Contact Dynamics,'' in the Oberwolfach Reports Open Problems collection. The record gives proposed year \(2013\) and identifies no individual proposer.
\item Supplied report citation: \emph{Dynamische Systeme (2014)}, Oberwolfach Reports, pp.\ 1975--2033, DOI \texttt{10.4171/owr/2013/34}. \url{https://doi.org/10.4171/owr/2013/34}. The background describes Oberwolfach Reports as records of mathematical workshops and their open-problem discussions. The literature-source list labels this reference ``Oberwolfach Report 34/2013.'' The differing year labels are retained as supplied.
\item Supplied literature reference: \emph{A dynamical characterization of universally tight lens spaces}. \url{https://arxiv.org/abs/1306.6617}. The record supplies no author list or precise theorem reference.
\item Supplied literature assessment, checked \texttt{2026-08-21}: status ``partially solved.'' The assessment reports a dynamical characterization for standard tight lens spaces, including the three-sphere, using a distinguished Reeb orbit and index/linking data, and states that no characterization of every closed connected three-manifold is known. This is supplied dated curation derived from the OWR material, not an independently verified literature assessment; time-sensitive status and cited claims require independent verification.
\item The registry labels the statement ``exact'' and supplies a verification note saying that its mathematical content was checked against the official Oberwolfach report, with notation, context, and question form normalized. These are provenance claims in the supplied record; no source was independently visited or checked for this draft.
\end{itemize}
\end{document}
